\section{Problem Formulation}
\label{sec:problem}

\subsection{Data roles}

A coordinator holds a public encoder $g$, an unlabelled \emph{proxy pool}
$P$ on which regions are formed, a \emph{development pool} whose labels are
revealed only under a declared acquisition mode, and a \emph{test pool}
used only offline. The class count and taxonomy are public; class
frequencies and external validation labels are not free.

Private data are held by $n$ clients. Client $i$ holds a labelled dataset
$D_i$ and may train a local model $f_i$; the coordinator observes clients
only through a protected sum over all participants, and the unit of
protection is a complete client dataset. For private verification, each
client also holds held-out records used only for bounded validation reports.

\subsection{Resources}

Each configuration carries a resource vector. \emph{Annotation}: $m$
distinct development labels and an acquisition mode, either \emph{random}
(uniform draws), \emph{balanced} ($m/C$ per class) or
\emph{feature-directed} (examples nearest each public region, requested
before any label is revealed); every fitting, selection and public
publication decision uses labels within $m$. \emph{Privacy}: one total
$(\varepsilon,\delta)$ budget under replace-one whole-client adjacency,
composed over all releases of a deployment. \emph{Payload}: integer
coordinates per client and release, before secret sharing.
\emph{Public computation}: decoding, assignment, selection and an optional
public student. Configurations differing in more than one resource are
reported as separate conditions.

\subsection{Release objects and endpoints}

A release object is the protected vector a client contributes: a clipped
parameter update, per-example proxy votes, or one bounded opinion per public
region. The published endpoint is a \emph{public classifier}; client models
are never published.

\subsection{Marginal value of private supervision}

Let $\mathcal H_{\mathrm{pub}}$ be a fixed family of public learners, each of
which maps a label set $L_m$ and the public features to a classifier. The
encoder is part of the resource vector: a zero-shot head on another
encoder, or one supervised by class names, changes the representation, not
$m$, and is compared as a separate condition. Let
$U(h)$ denote test accuracy (balanced accuracy for Usenet).

\begin{definition}[Resource-matched marginal value]
\label{def:value}
For acquisition mode $a$, budget $m$ and total privacy budget $\varepsilon$,
\begin{equation}
V_a(m,\varepsilon)=\mathbb E\bigl[U(h_{\mathrm{joint}})\bigr]
-\max_{h\in\mathcal H_{\mathrm{pub}}}\mathbb E\bigl[U(h)\bigr],
\label{eq:marginal-value}
\end{equation}
where $h_{\mathrm{joint}}$ combines the same label set with one protected
release under budget $\varepsilon$, and expectations run over label draws
from mode $a$ and over cohort/noise seeds.
\end{definition}

No within-budget selection rule is guaranteed to reach the maximum over
learner \emph{means}, so Eq.~\eqref{eq:marginal-value} is conservative. It
depends on $a$ as well as $m$, and since a deployment sees one label draw and
one cohort, we also report the per-draw gain distribution.

\subsection{Publication targets}

A positive $V_a$ concerns average utility; a publishing coordinator may
require more. We consider two targets for a candidate $f$ against a fixed public comparator $b$. The \emph{gain target}
requires the finite held-out cohort mean $G=n^{-1}\sum_i
[u_i(f)-u_i(b)]$ to be positive, where $u_i$ is held-out accuracy at client
$i$. The \emph{classwise target} additionally requires every client-weighted
class recall $\sum_i \mathrm{TP}_{i,c}/H_i\,/\sum_i N_{i,c}/H_i$ to be positive
($H_i$ held-out records at client $i$, $N_{i,c}$ of class $c$); it is the
pooled recall when all $H_i$ are equal, as in the CIFAR studies. Evidence can come from public labels or from protected private statistics
charged to the same total budget.


\paragraph{Research questions}
They follow the platform's three decisions. \textbf{Q1} (buy): how does
$V_a(m,\varepsilon)$ depend on acquisition mode and budget, and how is the
per-draw gain distributed? \textbf{Q2} (release): does report sensitivity rank
regional, per-example and head releases? \textbf{Q3} (verify): what does each
publication target cost in total budget, and can the budget split be chosen
before deployment? We then ask where, below the partition ceiling, regional
value is lost. Sections~\ref{sec:eval:main}, \ref{sec:eval:rounds},
\ref{sec:eval:verification} and~\ref{sec:eval:ceiling} answer them in turn.
